\documentclass[autodetect-engine,dvipdfmx,platex]{deimj}
\usepackage{latexsym}
\usepackage{booktabs}
\usepackage{color}
\usepackage{colortbl}
\usepackage[table]{xcolor}
\usepackage{graphicx}
\usepackage{amsmath}
\usepackage{bm}
\usepackage{multirow}
\usepackage{hyperref}
\usepackage{pxjahyper}
\usepackage[switch]{lineno}
\usepackage{tabularx}
\usepackage{makecell}
\usepackage{arydshln}  % 点線（dashed line）用
% \usepackage{placeins}



% 印刷位置調整 %
% 必要に応じて値を変更してください．
\hoffset -10mm % <-- 左に 10mm 移動
\voffset -10mm % <-- 上に 10mm 移動

\newcommand{\AmSLaTeX}{%
 $\mathcal A$\lower.4ex\hbox{$\!\mathcal M\!$}$\mathcal S$-\LaTeX}
\newcommand{\PS}{{\scshape Post\-Script}}
\def\BibTeX{{\rmfamily B\kern-.05em{\scshape i\kern-.025em b}\kern-.08em
 T\kern-.1667em\lower.7ex\hbox{E}\kern-.125em X}}

% \papernumber{DEIM Forum 2026 XX-Y}

\jtitle{複数基準特徴選択とOOB重み付けアンサンブルを用いた欠損値補完手法}
\authorlist{%
 \authorentry[c5625282@aoyama.jp]{高良 流平}{Ryuhei Kora}{AogakuUniv}%
 \authorentry[saitoh@ise.aoyama.ac.jp]{齊藤 史哲}{Fumiaki Saito}{AogakuUniv}%
}
\affiliate[AogakuUniv]{青山学院大学理工学研究科\hskip1zw
  〒252--5258 神奈川県相模原市中央区淵野辺本町 5--10--1}
 {Faculty of Informatics, Aoyamagakuin University\\
  5--10--1, Fuchinobe-Honcho, Chuo-ku, Sagamihara-shi, Kanagawa, 252--5258, Japan}

\begin{document}
% \linenumbers %このオプションは投稿時には無効にすること
\pagestyle{empty}

%%%%%%%%%%%%%%%%%%%%%%%%%%%%
% アブストラクト（300字程度）
%%%%%%%%%%%%%%%%%%%%%%%%%%%%
\begin{jabstract}
現実のデータ収集環境では，センサーの途切れや入力漏れ，システム障害などにより欠損値が頻繁に発生する．欠損値が存在すると，統計解析や機械学習モデルの推定に歪みが生じ，分析結果の信頼性が低下するため，適切な欠損処理が不可欠である．欠損処理には削除と補完があるが，削除は情報損失を伴う．一方，平均値補完などの統計的手法は，データの非線形性や特徴量間の関係を十分に捉えられない．この課題に対し，ランダムフォレストを用いた missForest が提案されているが，すべての特徴量を一律に用いるため，寄与の小さい特徴量の影響を受けやすい．そこで本研究では，複数の特徴選択基準に基づく特徴量部分集合を用い，OOB 指標に基づいて予測を統合するアンサンブル型欠損値補完手法を提案する．

\end{jabstract}

%%%%%%%%%%%%%%%%%%%%%%%%%%%%
% キーワード
%%%%%%%%%%%%%%%%%%%%%%%%%%%%
\begin{jkeyword}
欠損値補完，ランダムフォレスト，アンサンブル学習，特徴選択
\vspace{2em}

\end{jkeyword}


%%%%%%%%%%%%%%%%%%%%%%%%%%%%
% 本文
%%%%%%%%%%%%%%%%%%%%%%%%%%%%
\maketitle

% 0章：LateXの使い方
% \input{../contents/text/latex}

% 1章：はじめに
\input{../contents/text/introduction}

% 2章：関連研究
\section{関連研究}
\subsection{MissForestの概要}
missForest は，ランダムフォレストを基盤とした反復型の欠損値補完手法であり，分布仮定を必要としない
非パラメトリックな手法として提案されている．各変数を目的変数，残りの変数を説明変数とみなして
ランダムフォレストを学習し，欠損部分を予測する操作を変数ごとに繰り返すことで，データ行列全体の欠損値を
補完する枠組みを採用している．この反復的な学習・更新により，特徴量間の非線形な関係や相互作用を考慮した
補完が可能である．

missForest の特徴として，連続変数とカテゴリ変数が混在するデータに適用可能である点や，補完精度の推定に
ランダムフォレストの OOB誤差を利用できる点が挙げられる．OOB 誤差を用いることで，
外部の検証用データを用いることなく，補完結果の誤差を近似的に評価できる．また，多数の決定木の予測結果を
平均化するランダムフォレストの構造は，多重代入的な性質を内包していると整理されている．

原著論文およびその後の研究において，missForest は医療データや生物データをはじめとする実データを用いた
実験により，他の欠損値補完手法と比較して安定した補完性能を示すことが報告されており，機械学習に基づく
欠損値補完手法の代表的な手法として位置づけられている．一方で，高次元データにおいては，補完に寄与しない
特徴量の影響を受けやすいことが指摘されており，特徴量の扱いに関して改良の余地があるとされている．


% 3章：提案事項
\input{../contents/text/approach}

% 4章：評価実験
\input{../contents/text/experiment}

% 5章：結果
% \input{../contents/text/results}

% 6章：考察
\input{../contents/text/discussion}

% 7章：おわりに
\section{おわりに}
% 本研究ではアンサンブルベースの欠損補完方法の代表的手法であるmissForestとの比較において特徴量選択を導入した．ベンチマークデータによる性能評価により改善が確認できた．

本研究では，アンサンブルベースの欠損補完手法の代表例である missForest を基盤とし，複数の特徴選択基準を導入した欠損値補完手法を提案した．提案手法では，補完対象となる特徴量ごとに有効な説明変数の部分集合を構築し，各部分集合に基づいて学習したランダムフォレストを OOB 指標により統合することで，特徴選択の多様性と補完性能の安定化を図った．
ベンチマークデータおよび OpenML データセットを用いた実験により，提案手法は missForest や kNN 補完と比較して，多くの条件下で同等以上の補完精度を示すことを確認した．特に，中程度以上の欠損率や特徴量間の相関構造を持つデータにおいて，特徴選択とアンサンブル統合の効果が有効に機能することが示唆された．
一方で，本研究では連続変数のみを対象としており，カテゴリカル変数を含むデータへの適用や，欠損メカニズムが MCAR 以外の場合の検討は今後の課題である．今後は，より複雑な欠損構造や実データへの適用を通じて，提案手法の汎用性と実用性について検証を進めていく予定である．


\vspace{2em}
% 謝辞
\input{deim-acknowledgment}


%\vspace{30mm} <- 文献が本文と近すぎるときは適宜利用してください．
\vspace{2em}

\begin{thebibliography}{99}
\bibitem{textbook}
  高橋 将宣，渡辺 美智子，欠測データ処理 Rによる単一代入法と多重代入法，共立出版，(2017).
\bibitem{MissForest}
  Stekhoven, D. J. and Buhlmann, P., MissForest: non-parametric missing value imputation for mixed-type data, Bioinformatics, Vol. 28, No. 1, pp. 112-118 (2012). 
\bibitem{medical}
Chen, Z., Tan, S., Chajewska, U., Rudin, C. and Caruana, R., Missing Values and Imputation in Healthcare Data: Can Interpretable Machine Learning Help?, Proceedings of Machine Learning Research, Vol. 209, pp. 86-99 (2023).
\bibitem{industry}
Ma, L., Wang, M. and Peng, K., A missing manufacturing process data imputation framework for nonlinear dynamic soft sensor modeling and its application, Expert Systems with Applications, Vol. 237, Part A, Article 121428 (2024).
\bibitem{Rubin1976}
  Rubin, D. B., Inference and Missing Data, Biometrika, Vol. 63, No. 3, pp. 581--592 (1976).

\bibitem{LittleRubin2002}
  Little, R. J. A. and Rubin, D. B., Statistical Analysis with Missing Data (2nd ed.), Wiley, New York, NY (2002).

\bibitem{SchaferGraham2002}
  Schafer, J. L. and Graham, J. W., Missing data: our view of the state of the art, Psychological Methods, Vol. 7, No. 2, pp. 147--177 (2002).

\bibitem{Abe2016}
  阿部 貴行，欠測データの統計解析，朝倉書店，(2016).

\bibitem{AndridgeLittle2010}
  Andridge, R. R. and Little, R. J. A., A Review of Hot Deck Imputation for Survey Non-response, International Statistical Review, Vol. 78, No. 1, pp. 40--64 (2010).

\bibitem{BrickKalton1996}
  Brick, J. M. and Kalton, G., Handling missing data in survey research, Statistical Methods in Medical Research, Vol. 5, No. 3, pp. 215--238 (1996).

\bibitem{Rubin1987}
  Rubin, D. B., Multiple Imputation for Nonresponse in Surveys, John Wiley \& Sons, New York, NY (1987).

\bibitem{Breiman2001}
  Breiman, L., Random Forests, Machine Learning, Vol. 45, No. 1, pp. 5--32 (2001).

\bibitem{Strobl2007}
  Strobl, C., Boulesteix, A.-L., Zeileis, A. and Hothorn, T., Bias in random forest variable importance measures, BMC Bioinformatics, Vol. 8, Article 25 (2007).

\bibitem{Breiman1996}
  Breiman, L., Bagging predictors, Machine Learning, Vol. 24, No. 2, pp. 123--140 (1996).

\bibitem{GuyonElisseeff2003}
  Guyon, I. and Elisseeff, A., An Introduction to Variable and Feature Selection, Journal of Machine Learning Research, Vol. 3, pp. 1157--1182 (2003).

\bibitem{Hall1999}
  Hall, M. A., Correlation-based Feature Selection for Machine Learning, Ph.D. thesis, University of Waikato, Hamilton, New Zealand (1999).

\bibitem{DiazUriarte2006}
  D{\'i}az-Uriarte, R. and Alvarez de Andr{\'e}s, S., Gene selection and classification of microarray data using random forest, BMC Bioinformatics, Vol. 7, Article 3 (2006).

\end{thebibliography}



\end{document}
